\hypertarget{template-metaprogramming-power-features}{%
\chapter{TEMPLATE METAPROGRAMMING \& POWER
FEATURES}\label{template-metaprogramming-power-features}}

\hypertarget{variadic-templates}{%
\section{1. Variadic Templates}\label{variadic-templates}}

Templates that accept an arbitrary number of parameters.

\begin{lstlisting}[language={C++}]
// Base case (recursion terminator)
void print() {}

// Recursive step
template<typename T, typename... Args>
void print(T first, Args... args) {
    std::cout << first << " ";
    print(args...); // Unpack
}

int main() {
    print(1, "hello", 3.14);
}
\end{lstlisting}

\hypertarget{parameter-packing-...}{%
\subsection{\texorpdfstring{1.1 Parameter Packing
(\texttt{...})}{1.1 Parameter Packing (...)}}\label{parameter-packing-...}}

\begin{itemize}
\tightlist
\item
  \passthrough{\lstinline!typename... Args!}: Template parameter pack.
\item
  \passthrough{\lstinline!Args... args!}: Function parameter pack.
\item
  \passthrough{\lstinline!sizeof...(Args)!}: Number of arguments.
\end{itemize}

\begin{center}\rule{0.5\linewidth}{0.5pt}\end{center}

\hypertarget{type-traits}{%
\section{2. Type Traits}\label{type-traits}}

Compile-time type inspection and modification.

\begin{lstlisting}[language={C++}]
#include <type_traits>

static_assert(std::is_integral<int>::value, "Must be int");
static_assert(std::is_pointer<int*>::value, "Must be ptr");

// Conditional compilation (SFINAE)
template<typename T>
typename std::enable_if<std::is_integral<T>::value>::type
process(T x) {
    // Only compiles for integers
}
\end{lstlisting}

\begin{center}\rule{0.5\linewidth}{0.5pt}\end{center}

\hypertarget{using-aliases}{%
\section{\texorpdfstring{3. \texttt{using}
Aliases}{3. using Aliases}}\label{using-aliases}}

Replaces \passthrough{\lstinline!typedef!}, works with templates.

\begin{lstlisting}[language={C++}]
template<typename T>
using Dictionary = std::map<std::string, T>;

Dictionary<int> scores;
\end{lstlisting}

\begin{center}\rule{0.5\linewidth}{0.5pt}\end{center}

\hypertarget{stdtuple}{%
\section{\texorpdfstring{4.
\texttt{std::tuple}}{4. std::tuple}}\label{stdtuple}}

Generalization of \passthrough{\lstinline!std::pair!} to N elements.

\begin{lstlisting}[language={C++}]
#include <tuple>

std::tuple<int, double, std::string> t(1, 3.14, "hi");
int i = std::get<0>(t);
\end{lstlisting}
